\chapter*{Apresentação}

Este manual orienta a elaboração de Trabalhos de Conclusão de Curso no âmbito da Universidade do Distrito Federal Professor Jorge Amaury Maia Nunes -- UnDF. Sua finalidade é apoiar estudantes, docentes, orientadores, membros de banca e coordenações de curso na construção de trabalhos acadêmicos claros, rigorosos, éticos e formalmente padronizados.

O documento combina orientação acadêmica, critérios gerais de escrita científica, organização metodológica e normalização. Ele não é um manual de software e não tem como foco ensinar ferramentas de edição. Quando menciona padrões visuais, elementos obrigatórios ou decisões de apresentação, a intenção é registrar o padrão institucional que deve orientar a versão final do trabalho.

A normalização adotada considera as normas ABNT vigentes utilizadas pelo projeto e as decisões institucionais validadas durante a construção do modelo de TCC da UnDF. Nos pontos em que a ABNT admite escolhas, o manual apresenta o padrão definido e incorporado aos artefatos do projeto.

\caixaatencao{Este manual não substitui regulamentos específicos de curso, deliberações de colegiado, normas de comitês de ética ou exigências formais publicadas pela instituição. Quando houver regra específica do curso, ela deve ser observada em conjunto com este manual.}

\section*{Objetivos do manual}

São objetivos deste manual:

\begin{itemize}
  \item apresentar o sentido acadêmico do TCC e seu papel na formação superior;
  \item orientar a escolha do tema, a definição do problema e a construção dos objetivos;
  \item explicar a estrutura esperada de trabalhos acadêmicos;
  \item diferenciar elementos obrigatórios, opcionais e condicionais;
  \item orientar a escrita acadêmica, a integridade autoral e o uso responsável de fontes;
  \item esclarecer critérios gerais de apresentação, normalização, entrega e avaliação;
  \item oferecer roteiros, checklists e instrumentos de apoio ao estudante.
\end{itemize}

\section*{Como ler este manual}

O estudante que está iniciando o TCC deve ler os capítulos sobre princípios, percurso, tema, problema e projeto. O estudante em fase de escrita deve priorizar os capítulos sobre estrutura, elementos textuais, escrita acadêmica, citações e referências. O estudante em fase de entrega deve consultar os capítulos sobre normalização, avaliação, integridade, entrega e checklists.

O orientador pode usar este manual como referência comum para alinhar expectativas com seus orientandos. A banca pode utilizá-lo como apoio para verificar se a forma do trabalho está coerente com o padrão institucional, sem confundir questões de forma com a avaliação do mérito acadêmico.
